\subsection{A sharp inverse estimate returned to the manuscript's original directions}
\label{sec:ns-reverse-smooth}

This subsection applies the preceding coordinate-level operator calculation
back to the public manuscript's equation~(6.7), equation~(8.10), and
Lemma~8.6, equation~(8.19). In the preserved 166-page extraction these
occur on printed pages~64, 91, and 95, respectively. The manuscript's
inverse acts on the continuous torus with probability Haar measure.
We retain that domain and the original directions
\[
 \begin{gathered}
 v_r=(1,1-\sqrt2),\qquad v_t=(\sqrt2-1,1),\\
 D_v=v_1\partial_{Y_1}+v_2\partial_{Y_2},\qquad
 Y\in\mathbb R^2/\mathbb Z^2.
 \end{gathered}
\]
The following estimates strengthen the constants and derivative count
in those inverse steps. They require no PDE existence assertion.

\begin{theorem}[Sharp denominator constant in the original coordinates]
\label{thm:ns-sharp-denominator}
Put
\[
 \begin{gathered}
 c=\sqrt{4+2\sqrt2},\qquad d=\sqrt{4-2\sqrt2},\\
 w_r=(1,1+\sqrt2),\qquad w_t=(-1-\sqrt2,1).
 \end{gathered}
\]
For either paired choice $(v,w)=(v_r,w_r)$ or $(v_t,w_t)$,
$v\cdot w=0$, $\|v\|_2=d$, $\|w\|_2=c$, and for every
$k=(m,n)\in\mathbb Z^2\setminus\{0\}$,
\[
 |v\cdot k|>\frac1{c\|k\|_2},\qquad
 \inf_{k\ne0}|v\cdot k|\|k\|_2=\frac1c.             \tag{S1}
\]
In particular the constant $1/c$ cannot be increased.
\end{theorem}
\begin{proof}
Direct multiplication gives all dot products and squared lengths.
The products of the original symbols with their conjugates are
\[
 \begin{split}
 (v_r\cdot k)(w_r\cdot k)&=m^2+2mn-n^2=(m+n)^2-2n^2,\\
 (v_t\cdot k)(w_t\cdot k)&=n^2-2mn-m^2=(n-m)^2-2m^2.
 \end{split}                                                     \tag{S2}
\]
Each is a nonzero integer: its vanishing with a nonzero second
coordinate would make $\sqrt2$ rational, and when that coordinate
vanishes the first one must also vanish. Irrationality follows from
the parity argument given in Theorem~\ref{thm:ns-directional-intertwining}.
Their absolute values are therefore at least one. Cauchy--Schwarz
gives $|w\cdot k|\leq c\|k\|_2$. Equality would make the nonzero
integer vector $k$ parallel to $w$, whose coordinate ratio is
irrational, and is impossible. Hence
$1\leq|v\cdot k||w\cdot k|<c|v\cdot k|\|k\|_2$.

To prove optimality, retain the explicit Pell sequence
\[
 \begin{gathered}
 \alpha=1+\sqrt2,\qquad\beta=1-\sqrt2=-\alpha^{-1},\\
 p_0=1,\quad q_0=0,\quad
 p_{j+1}=p_j+2q_j,\quad q_{j+1}=p_j+q_j.
 \end{gathered}                                                   \tag{S3}
\]
Induction proves $p_j,q_j$ are integers and
$p_j+q_j\sqrt2=\alpha^j$, $p_j-q_j\sqrt2=\beta^j$.
In particular $p_j^2-2q_j^2=(-1)^j$. For $j\geq1$ set
\[
 k_{r,j}=(p_j-q_j,q_j),\qquad
 k_{t,j}=(q_j,q_j-p_j).
\]
The two vectors have the same length $K_j$. Their actual symbol
and conjugate values are respectively
\[
 \begin{array}{c|cc}
 &v\cdot k_{v,j}&w\cdot k_{v,j}\\\hline
 r&\beta^j&\alpha^j\\
 t&-\beta^j&-\alpha^j
 \end{array}
\]
by substitution. Since the retained nonzero perpendicular vectors
$v,w$ span $\mathbb R^2$, the exact decomposition is
$k=(v\cdot k)v/d^2+(w\cdot k)w/c^2$. Therefore
\[
 K_j^2=\frac{\alpha^{2j}}{c^2}+\frac{\alpha^{-2j}}{d^2},
 \qquad
 |v\cdot k_{v,j}|^2K_j^2
       =\frac1{c^2}+\frac{\alpha^{-4j}}{d^2}.          \tag{S4}
\]
The number $\alpha$ is greater than one, so $K_j$ tends to infinity
and the second expression tends to $1/c^2$. This proves (S1)'s
infimum and its nonattainment with every integer coordinate retained.
\end{proof}

We specify the analytic spaces used below. For $Y\in[0,1]^2$ write
\[
 \begin{gathered}
 e_k(Y)=e^{2\pi i k\cdot Y},\qquad W(k)=1+4\pi^2\|k\|_2^2,\\
 \widehat f(k)=\int_{[0,1]^2}f(Y)e^{-2\pi i k\cdot Y}\,dY,
 \end{gathered}
\]
For real $s$, $H^s$ is the completion of finite trigonometric
polynomials with the exact norm
\[
 \|f\|_{H^s}^2=\sum_{k\in\mathbb Z^2}W(k)^s|\widehat f(k)|^2,
 \qquad H^s_0=\{f\in H^s:\widehat f(0)=0\}.           \tag{S5}
\]
Thus elements may be represented by their weighted square-summable
coefficient arrays; truncation of such arrays converges in this norm
because the defining nonnegative series has vanishing tails.
For integer $m\geq0$ the coordinate norm on $C^m$ is
\[
 \|f\|_{C^m}=\max_{a,b\geq0,\ a+b\leq m}
          \sup_{Y\in\mathbb T^2}|\partial_{Y_1}^a\partial_{Y_2}^bf(Y)|.
                                                               \tag{S6}
\]
All integrations and suprema use the original period-one coordinates.

\begin{lemma}[Fourier reconstruction and the precise analytic estimates used]
\label{lem:ns-fourier-analytic-bridge}
For every smooth periodic function $f$, its Fourier series and each
coordinatewise differentiated series converge absolutely and uniformly
to $f$ and its corresponding derivative. For an integer $r\geq0$,
\[
 \|f\|_{H^r}\leq3^{r/2}\|f\|_{C^r}.                \tag{S7}
\]
The positive lattice constant
\[
 S=\sum_{k\in\mathbb Z^2\setminus\{0\}}W(k)^{-2}
       \quad\hbox{satisfies}\quad S\leq\frac3{4\pi^4}.            \tag{S8}
\]
For every zero-mean coefficient array $b$ and integer $m\geq0$,
\[
 \sum_{k\ne0}(2\pi\|k\|_2)^m|b(k)|
    \leq\sqrt S\left(\sum_{k\ne0}W(k)^{m+2}|b(k)|^2\right)^{1/2}.
                                                               \tag{S9}
\]
Whenever the right side is finite, the corresponding Fourier series
and its derivatives through order $m$ converge absolutely and
uniformly, with $C^m$ norm at most that right side.
\end{lemma}
\begin{proof}
The functions $e_k$ are orthonormal, since integrating each coordinate
character gives one at integer frequency zero and zero elsewhere.
For a finite set $F$ of frequencies, expansion of
$\|g-\sum_{k\in F}\widehat g(k)e_k\|_{L^2}^2\geq0$
therefore gives
$\sum_{k\in F}|\widehat g(k)|^2\leq\int|g|^2$.
Increasing $F$ proves Bessel's inequality for every continuous $g$.
Periodic integration by parts gives
$\widehat{\partial_1^a\partial_2^bf}(k)
 =(2\pi i k_1)^a(2\pi i k_2)^b\widehat f(k)$;
the boundary terms cancel at the two equal endpoint values.
Expanding $W(k)^r=(1+(2\pi k_1)^2+(2\pi k_2)^2)^r$
by the multinomial formula and applying Bessel to each derivative
with orders $a,b$, $a+b\leq r$, gives
\[
 \begin{split}
 \sum_kW(k)^r|\widehat f(k)|^2
 &\leq\sum_{j+a+b=r}\frac{r!}{j!a!b!}
                   \int_{[0,1]^2}|\partial_1^a\partial_2^bf|^2\\
 &\leq3^r\|f\|_{C^r}^2.
 \end{split}
\]
This proves (S7) without requiring an unproved Parseval identity.

For a nonzero frequency choose a coordinate with
$|k_i|\geq\|k\|_2/\sqrt2$. Integrating by parts $M$ times in
that coordinate yields
\[
 |\widehat f(k)|\leq
       \left(\frac{\sqrt2}{2\pi\|k\|_2}\right)^M\|f\|_{C^M}.
                                                               \tag{S10}
\]
Exactly $8q$ integer points have $\max(|k_1|,|k_2|)=q$.
Thus $\sum_{k\ne0}\|k\|_2^{-t}\leq8\sum_{q\geq1}q^{1-t}$
is finite for $t>2$. Formula (S10) with arbitrarily large $M$
proves absolute uniform convergence of the series and every
differentiated series. Its continuous sum has Fourier coefficients
$\widehat f(k)$ by uniform convergence and orthogonality.
The difference from $f$ has all coefficients zero. Its product
Fej\'er convolutions are zero and converge uniformly to that
difference by the explicit positive-kernel argument following
Theorem~\ref{thm:ns-torus-average}; hence the difference is zero.
This proves reconstruction and differentiation. The same argument
also yields Parseval for smooth functions if needed: finite Fourier
sums converge in $L^2$, and their squared norms are the finite sums
of squared coefficients by orthogonality.

The same $8q$ count gives
\[
 S\leq\sum_{q=1}^{\infty}\frac{8q}{(1+4\pi^2q^2)^2}
 \leq\frac1{2\pi^4}\sum_{q=1}^{\infty}q^{-3}
 \leq\frac1{2\pi^4}\left(1+\int_1^{\infty}t^{-3}\,dt\right)
 =\frac3{4\pi^4}.
\]
Finally $(2\pi\|k\|_2)^m\leq W(k)^{m/2}$.
Apply Cauchy--Schwarz to
$W(k)^{-1}$ and $W(k)^{(m+2)/2}|b(k)|$ to prove (S9).
For a derivative of order $a+b\leq m$, its multiplier is bounded
by $(2\pi\|k\|_2)^{a+b}\leq(2\pi\|k\|_2)^m$ because
$\|k\|_2\geq1$. The resulting summable majorant is uniform in
$Y$, proving all stated convergence and $C^m$ assertions.
\end{proof}

\begin{theorem}[The smooth inverse, exact Sobolev norm, and sharp derivative loss]
\label{thm:ns-sharp-smooth-inverse}
For either original direction $v$ and an integer $q\geq1$, define
\[
 \widehat{T_v^q f}(0)=0,\qquad
 \widehat{T_v^q f}(k)
       =\frac{\widehat f(k)}{(2\pi i\,v\cdot k)^q}\quad(k\ne0).
                                                               \tag{S11}
\]
For every real $s$, this is a bounded map
$T_v^q:H^{s+q}_0\to H^s_0$ with exact operator norm
\[
 \|T_v^q\|_{H^{s+q}_0\to H^s_0}
       =\left(\frac c{4\pi^2}\right)^q.              \tag{S12}
\]
For every $\delta<q$, this inverse has no bounded extension
$H^{s+\delta}_0\to H^s_0$ agreeing with (S11) on polynomials.

For smooth zero-mean $f$, $T_v^qf$ is the unique smooth zero-mean
solution of $D_v^q u=f$. More generally, for every datum
$f\in H^{s+q}_0$, its coefficients give the unique zero-mean
distributional solution. Here a distribution is a continuous
complex-linear functional on $C^\infty(\mathbb T^2)$ with its
seminorms (S6), and its coefficient at $k$ is its value on $e_{-k}$.
On smooth periodic
functions the kernel of $D_v^q$ is precisely the constants, its
image is all smooth zero-mean functions, and its full fibre at $f$
is $T_v^q f+\mathbb C$.
\end{theorem}
\begin{proof}
Theorem~\ref{thm:ns-sharp-denominator} gives the exact pointwise bound
\[
 \frac1{|2\pi i\,v\cdot k|^q W(k)^{q/2}}
 <\left(\frac{c\|k\|_2}{2\pi\sqrt{1+4\pi^2\|k\|_2^2}}\right)^q
 <\left(\frac c{4\pi^2}\right)^q.                   \tag{S13}
\]
Multiplying coefficient by coefficient and summing nonnegative
squares proves the upper bound in (S12) on polynomials and all
weighted square-summable arrays. Thus truncation extends the same
formula continuously to the stated completion, with no change of
coordinates. On the single character at $k_{v,j}$ of (S3)--(S4),
the ratio of the output norm to the input norm is
\[
 \frac{\alpha^{qj}}{(2\pi)^q(1+4\pi^2K_j^2)^{q/2}}
        \longrightarrow\left(\frac c{4\pi^2}\right)^q,
\]
because $K_j/\alpha^j\to1/c$ by (S4). This proves the exact
operator norm. If the input order is $s+\delta$ instead, the ratio
is
$\alpha^{qj}/((2\pi)^q(1+4\pi^2K_j^2)^{\delta/2})$.
Its quotient by $\alpha^{(q-\delta)j}$ tends to the positive
constant $c^\delta/(2\pi)^{q+\delta}$. It is unbounded when
$\delta<q$, proving sharpness of the derivative loss.

For smooth $f$, (S10) and the multiplier growth
$|(2\pi i\,v\cdot k)^{-q}|\leq(c\|k\|_2/(2\pi))^q$
give uniform absolute convergence of every differentiated inverse
series. Lemma~\ref{lem:ns-fourier-analytic-bridge} and termwise
differentiation therefore show smoothness and $D_v^qT_v^qf=f$.
The zero coefficient is exactly its Haar mean. Every nonzero
Fourier multiplier of $D_v^q$ is nonzero, so any smooth solution
has exactly the coefficients (S11) away from zero. Its remaining
constant is arbitrary; Fourier reconstruction proves the full
kernel and fibres.

To make the distributional assertion explicit, the resulting
$u\in H^s_0$ has coefficients of at most polynomial growth:
each term of its squared norm is bounded by that full norm, so
$|\widehat u(k)|\leq\|u\|_{H^s}W(k)^{-s/2}$.
Consequently there are $B\geq0$ and an integer $r\geq0$ with
$|\widehat u(k)|\leq B\|k\|_2^r$ for $k\ne0$;
one obtains them directly from
$W(k)\leq(1+4\pi^2)\|k\|_2^2$ and $W(k)\geq1$.
Define its pairing with a smooth periodic test function by
$u(\varphi)=\sum_k\widehat u(k)\widehat\varphi(-k)$.
Choose an integer $M>r+2$. Formula (S10) and the lattice count
give the explicit continuous-functional bound
\[
 |u(\varphi)|\leq
 \left(|\widehat u(0)|+
    8B\left(\frac{\sqrt2}{2\pi}\right)^M
       \sum_{j=1}^{\infty}j^{1+r-M}\right)\|\varphi\|_{C^M}.
\]
The series converges, so this defines a distribution with exactly
the specified coefficients. The same construction realizes
$f\in H^{s+q}_0$ as a distribution. Distributional differentiation
is defined by $(D_vu)(\varphi)=-u(D_v\varphi)$; evaluating on
$e_{-k}$ therefore multiplies its coefficient by $2\pi i\,v\cdot k$.
Formula (S11) gives the coefficient equation for $D_v^qu=f$.
Finally a distribution with all Fourier coefficients zero vanishes
on every trigonometric polynomial. The Fourier partial sums of
any smooth test function converge to it in every seminorm (S6),
by Lemma~\ref{lem:ns-fourier-analytic-bridge}. Continuity of the
distribution therefore makes its value on every smooth test
function zero. This proves the coefficient equation as an equality
of distributions and proves uniqueness: any other zero-mean
solution has the same forced nonzero coefficients and the same
zero coefficient.
\end{proof}

\begin{theorem}[A strengthened coordinate derivative estimate]
\label{thm:ns-cm-improvement}
For integers $m\geq0$, $q\geq1$ and zero-mean
$f\in C^{m+q+2}(\mathbb T^2)$, the inverse series (S11) belongs
to $C^m$ and satisfies
\[
 \begin{aligned}
 \|T_v^qf\|_{C^m}
   &\leq\left(\frac c{4\pi^2}\right)^q
             \sqrt S\,\|f\|_{H^{m+q+2}}\\
   &\leq C_{m,q}\|f\|_{C^{m+q+2}}.
 \end{aligned}                                                   \tag{S14}
\]
where the completely specified constant is
\[
 C_{m,q}=\left(\frac c{4\pi^2}\right)^q
          3^{(m+q+2)/2}\sqrt S
 \leq\left(\frac c{4\pi^2}\right)^q
          \frac{3^{(m+q+3)/2}}{2\pi^2}.             \tag{S15}
\]
For $q=1$ this uses $m+3$ input derivatives and applies to the
original temporal inverse in equation~(8.19). For general $q$ it
uses $m+q+2$ and applies to the original repeated radial inverse
in equation~(8.10). No optimality assertion is made for this
particular $C^m$ derivative count.
\end{theorem}
\begin{proof}
Apply (S9) to the inverse coefficients, and then (S13) with
$s=m+2$. This gives the first inequality (S14), including uniform
convergence of all inverse derivatives through order $m$.
Formula (S7) with $r=m+q+2$ gives the second inequality, and
(S8) gives the displayed bound on the constant. The proofs of
(S7) and (S9) only require the finite number of continuous
derivatives displayed here, so the same argument applies directly
to $C^{m+q+2}$ data. In the two source equations the operators
are exactly $T_{v_t}$ and $T_{v_r}^q$, with the same Fourier
constant $2\pi i$, sign, period and zero mode. The source bounds
with $m+4$ and $m+q+3$ input derivatives follow at once from
(S14), since increasing $C^r$ order only increases the norm (S6).
\end{proof}

\begin{proposition}[Parameter derivatives, support, reality and the source mean update]
\label{prop:ns-inverse-parameters}
Let $V\subset\mathbb R^b$ be open, and let
$f:V\times\mathbb T^2\to\mathbb C$ be jointly smooth with
$\int_{\mathbb T^2}f(x,Y)\,dY=0$ for every $x\in V$.
Define $u(x,Y)=T_v^q(f(x,\cdot))(Y)$. Then $u$ is jointly
smooth and, for every parameter multi-index $\gamma$,
\[
 \begin{gathered}
 \partial_x^\gamma u=T_v^q(\partial_x^\gamma f),\\
 \sup_{x\in K}\|\partial_x^\gamma u(x,\cdot)\|_{C^m}
 \leq C_{m,q}\sup_{x\in K}
                \|\partial_x^\gamma f(x,\cdot)\|_{C^{m+q+2}}.
 \end{gathered}                                                  \tag{S16}
\]
for every compact $K\subset V$. If $f(x,Y)=0$ for all $Y$
whenever $x\notin A\subset V$, the same statement holds for $u$.
Real-valued data give a real-valued inverse.

In the manuscript's fixed local band/common-torus chart, retain
$Q>0$, $h=h_{\rm NS}$, $c_{i_0}=T_g^{i_0}Q^{1+h}>0$,
$T_g=4+\sqrt2$, and the distinct coordinates
\[
 \begin{gathered}
 Y\in\mathbb T^2_{\rm abs},\qquad y=J^{i_0}Y\in\mathbb T^2_{\rm com},
 \qquad P_{i_0}F(Y)=F(J^{i_0}Y),\\
 D_v^Y=v\cdot\partial_Y,\qquad D_v^y=v\cdot\partial_y,\qquad
 T_v^Y=(D_v^Y)^{-1},\qquad T_v^y=(D_v^y)^{-1}.
 \end{gathered}
\]
The inverse in each coordinate acts on that torus's smooth
zero-Haar-mean functions and uses its own original Fourier
frequency with the multiplier (S11). Denote the corresponding
means by $\Pi_Y$ and $\Pi_y$. On the indicated smooth domains,
\[
 \begin{gathered}
 \Pi_Y P_{i_0}=\Pi_y,\qquad
 D_{v_t}^Y P_{i_0}=T_g^{i_0}P_{i_0}D_{v_t}^y,\\
 T_{v_t}^Y P_{i_0}=T_g^{-i_0}P_{i_0}T_{v_t}^y
       \quad\hbox{on }\ker\Pi_y.
 \end{gathered}
\]
For the actual common-chart residual write
$E_\theta^\circ=E_\theta-\Pi_yE_\theta$. Its temporal update is
\[
 \begin{gathered}
 \Delta v=-c_{i_0}^{-1}T_{v_t}^yE_\theta^\circ,\qquad
 c_{i_0}D_{v_t}^y\Delta v=-E_\theta^\circ,\\
 P_{i_0}\Delta v=-Q^{-1-h}T_{v_t}^Y(P_{i_0}E_\theta^\circ),\\
 \|\Delta v\|_{C_y^m}
       \leq c_{i_0}^{-1}C_{m,1}\|E_\theta^\circ\|_{C_y^{m+3}}.
 \end{gathered}
                                                               \tag{S17}
\]
The identical typed formulas apply to the manuscript's desired axial
increment $\gamma_d=-c_{i_0}^{-1}T_{v_t}^yE_z^\circ$.
For smooth shell-supported data the source's actual axial increment is
$\Delta\gamma=\gamma_d+a$, with $a=-A_1\gamma_d$, where its
retained cutoff operator is explicitly
\[
 \begin{split}
 (A_1f)(R,Z,T,y)
  ={}&R^{-1}(\partial_R\chi_m)(R,Z,T)\\
  &{}\cdot\int_0^\infty R'f\bigl(R',Z,T,
          y+M((R')^{d_r}-R^{d_r})v_r\bigr)\,dR',
 \qquad M=\Lambda^{i_0}Q^{d_r/2}.
 \end{split}
\]
Here $R>0$, the smooth zero extension of $f$ has radial support
locally bounded in a compact subinterval of $(0,\infty)$, and
$\chi_m(R,Z,T)$ is the source's smooth torus-independent cutoff.
The actual fast-time cancellation is consequently
\[
 c_{i_0}D_{v_t}^y\Delta\gamma
       =-E_z^\circ+c_{i_0}D_{v_t}^y a.
\]
The complete displayed time operator remains
$-\varepsilon\partial_T+c_{i_0}D_{v_t}^y$, with
$\varepsilon=Q^h$ and $\partial_T$ holding $y$ fixed. Thus its
action on the updates still contains
$-\varepsilon\partial_T\Delta v$ and
$-\varepsilon\partial_T\Delta\gamma$, respectively.
No zero or flatness assertion about $a$ is part of this inverse estimate.
\end{proposition}
\begin{proof}
On a compact parameter neighborhood inside $V$, every mixed
parameter/torus derivative of $f$ is bounded. Differentiation of
its coefficient integral in a parameter is justified by the
continuous bounded derivative on that compact set times the
finite torus measure. Repetition gives
$\partial_x^\gamma\widehat f(x,k)
=\widehat{\partial_x^\gamma f}(x,k)$ and also zero mean for
each parameter derivative. Integration by parts as in (S10)
gives every polynomial decay order uniformly on this neighborhood,
for all parameter derivatives under consideration. After multiplication
by the inverse multiplier and any fixed coordinate derivative,
choose the decay order larger than the total growth exponent plus
two. The $8q$ lattice-shell count then supplies a summable majorant
independent of the parameter and torus variables. Thus every mixed
series differentiates termwise and converges locally uniformly.
One may justify each first derivative directly by integrating its
uniformly convergent derivative series along a coordinate segment;
the fundamental theorem of calculus passes to the sum. Iterating
proves joint smoothness and (S16)'s exact commutation. Applying
(S14) at each parameter and taking suprema proves its bound.

When a parameter slice of $f$ vanishes, all its Fourier coefficients
vanish, and (S11) gives the vanishing slice of $u$. This proves the
support assertion without changing the parameter set. If $f$ is
real, $\widehat f(-k)=\overline{\widehat f(k)}$ and
$(2\pi i\,v\cdot(-k))^{-q}
=\overline{(2\pi i\,v\cdot k)^{-q}}$; pairing opposite modes
shows $u$ is real. A common-torus character $e_k(y)$ pulls back
to $e_{(J^{i_0})^Tk}(Y)$. Its absolute derivative multiplier is
$2\pi i\,v_t\cdot(J^{i_0})^Tk
=T_g^{i_0}2\pi i\,v_t\cdot k$.
The lattice map is injective and preserves zero frequency exactly,
which proves the stated Haar identity and both typed intertwining
identities, including the reciprocal factor for the inverse. Smooth
Fourier convergence proved above passes them from characters to the
specified smooth spaces. In particular
$P_{i_0}T_{v_t}^y=T_g^{i_0}T_{v_t}^YP_{i_0}$ on $\ker\Pi_y$.
Multiplying this identity by
$-c_{i_0}^{-1}=-T_g^{-i_0}Q^{-1-h}$ gives the absolute
representative in (S17) with its complete coefficient.

The equation $D_{v_t}^yT_{v_t}^yE_\theta^\circ=E_\theta^\circ$
proves (S17)'s common-chart cancellation, including its minus sign.
The estimate is (S14), now in the explicitly named common
coordinates, multiplied by $c_{i_0}^{-1}$. The same calculation
with $E_z^\circ$ gives the desired $\gamma_d$ and its bound.
For the stated shell-supported inputs, differentiation under the
displayed $A_1$ integral is justified locally by a fixed compact
radial interval and bounded smooth derivatives on that interval;
$R$ stays positive. This proves smoothness of $a$ there.
Linearity of $D_{v_t}^y$ applied to
$\Delta\gamma=\gamma_d+a$ proves the actual fast cancellation
with the retained term $c_{i_0}D_{v_t}^ya$. Finally applying the
complete time operator to each update gives exactly the displayed
slow derivatives in addition to its fast part. No slow derivative
has been commuted through the radial cutoff operator.
\end{proof}

The static source dependency for this result is explicit.
\texttt{DiophantineGraph.lean} proves the algebraic norm lower bounds
at lines~94 and~104 and the reciprocal bounds with constant~6 at
lines~226 and~241. \texttt{TorusInverse.lean} retains
$\omega=2\pi i$ at line~52, defines the same original vectors at
line~195, and proves the inverse coefficient bound and smooth
inverse at lines~278 and~323. Its Haar measure at line~342 is the
product of two actual circle Haar probability measures.
\texttt{SmoothFourierData.lean}, lines~351 and~561, proves the
rapid-coefficient and reconstruction bridge for smooth periodic
functions; its line~565 returns the directional equation for
those actual functions. Thus rapid coefficients are not being
inserted here as an unproved surrogate for a smooth function.
Equations (S1), (S12), and (S14) provide sharper analytical constants
and derivative estimates for that same source chain. These are
independent mathematical proofs with a bounded static source audit;
no Lean process or complete PDE verification is asserted.
